% !TeX encoding = UTF-8
\documentclass[variant=guiaresuelta,cover=institutional,mode=screen]{ua-engineering-v6-article}
% !TeX encoding = UTF-8
\UASetUniversity{Universidad Austral}
\UASetFaculty{Facultad de Ingeniería}
\UASetDepartment{Departamento de Computación}
\UASetCampus{Pilar}
\UASetCareer{Ingeniería Informática}
\UASetCommission{Comisión A}
\UASetProfessor{Equipo docente}
\UASetSemester{Segundo cuatrimestre 2026}
\UASetCourse{Sistemas Distribuidos}
\UASetDate{2026}

\UASetClassNumber{08}
\UASetClassTitle{Causalidad, consenso y consistencia distribuida}
\title{Clase 08 --- Guía resuelta}
\author{\UAProfessor}
\date{\UADate}
\begin{document}
\maketitle
\section{Cómo leer las soluciones}
Son modelos de argumentación. Otras respuestas son válidas si declaran supuestos, preservan el invariante, hacen explícito el límite y proponen evidencia verificable.

\section{Criterio de autocorrección}
Después de resolver cada ejercicio, asigne de cero a dos puntos en cada dimensión:
\begin{itemize}
\item \textbf{Ejecución:} identifica actores, mensajes, estado local y orden de eventos.
\item \textbf{Garantía:} formula una propiedad observable o un invariante, no una etiqueta vaga.
\item \textbf{Falla:} ubica exactamente timeout, crash, pausa o partición y explica qué información se pierde.
\item \textbf{Mecanismo:} aplica correctamente la regla de reloj, mayoría, quórum, versión o espera.
\item \textbf{Costo y límite:} indica qué no resuelve el mecanismo y dónde paga latencia, rechazo, espacio o reparación.
\item \textbf{Evidencia:} propone un timeline, métrica o experimento capaz de refutar la respuesta.
\end{itemize}
Una solución de 10--12 puntos está lograda; 7--9 puntos necesita explicitar supuestos o costos; menos de 7 puntos suele indicar que se recordó un nombre sin reconstruir la ejecución. La autocorrección no exige que su redacción coincida con la respuesta orientativa: exige que otra persona pueda mover la falla y predecir el resultado sin ambigüedad.

\section{Ejercicios y respuestas}
\begin{uaexercise}
\textbf{1. Timeout e indistinguibilidad.} Buenos Aires envía una reserva a Madrid y recibe timeout. Enumere cuatro historias compatibles con esa observación. Para cada una, indique qué evidencia adicional permitiría distinguirla.
\end{uaexercise}
\begin{uaanswer}
El request puede no haber salido, perderse en la red, llegar a un proceso pausado, ejecutarse con respuesta perdida o quedar bloqueado por una partición. IDs idempotentes, trazas correlacionadas, log durable, estado consultable y una respuesta posterior permiten reunir evidencia. El timeout por sí solo solo describe la espera local.

\par\smallskip\textbf{Criterio de autocorrección.} La respuesta está completa si separa falta de respuesta de falta de ejecución y nombra una fuente de evidencia durable o correlacionada.
\end{uaanswer}

\begin{uaexercise}
\textbf{2. Happened-before.} Dibuje tres líneas de proceso ---BA, líder y Madrid--- para click, send, receive y append. Marque dos relaciones causales, una relación transitiva y un par concurrente.
\end{uaexercise}
\begin{uaanswer}
Una cadena válida es $a_1\rightarrow a_2\rightarrow l_1\rightarrow l_2$ y $m_1\rightarrow m_2\rightarrow l_3$. Antes de que el líder reciba ambas ramas, $a_2\parallel m_2$. La transitividad permite inferir $a_1\rightarrow l_2$.

\par\smallskip\textbf{Criterio de autocorrección.} Debe aparecer una flecha de envío--recepción, una inferencia transitiva y un par sin camino causal en ninguno de los sentidos.
\end{uaanswer}

\begin{uaexercise}
\textbf{3. Lamport.} Asigne timestamps de Lamport a la ejecución anterior. Justifique la regla $L\leftarrow\max(L,t_m)+1$ y construya un par con $L(a)<L(b)$ pero sin $a\rightarrow b$.
\end{uaexercise}
\begin{uaanswer}
Una asignación posible: BA 1,2; Madrid 1,2; líder 3,4,5. $\max$ conserva ambos pasados conocidos y $+1$ coloca la recepción después del envío. Eventos independientes pueden recibir 1 y 2 sin que exista mensaje causal entre ellos.

\par\smallskip\textbf{Criterio de autocorrección.} Debe justificar tanto la preservación del máximo conocido como el incremento local; un simple listado de números no alcanza.
\end{uaanswer}

\begin{uaexercise}
\textbf{4. Relojes vectoriales.} Con componentes [BA, Madrid, Líder], calcule los vectores hasta que el líder recibe ambas reservas. Explique por qué el resultado del líder es [2,2,3] y no [2,2,5].
\end{uaexercise}
\begin{uaanswer}
BA produce [1,0,0] y [2,0,0]; Madrid [0,1,0] y [0,2,0]; líder [2,0,1], [2,0,2], [2,2,3]. El tercer componente cuenta tres eventos locales. El 5 es una etiqueta escalar Lamport, no un contador de eventos del actor líder.

\par\smallskip\textbf{Criterio de autocorrección.} Debe calcular máximos componente a componente y distinguir el contador local del líder de la etiqueta escalar de Lamport.
\end{uaanswer}

\begin{uaexercise}
\textbf{5. Máquina de estados.} Defina estados y transiciones para VIP-01. Incluya reserva de Sofía, reserva de Tomás y repetición idempotente de un comando ya aplicado.
\end{uaexercise}
\begin{uaanswer}
Estados mínimos: \texttt{AVAILABLE}, \texttt{RESERVED\_BY\_SOFIA} y \texttt{RESERVED\_BY\_TOMAS}. Desde AVAILABLE cada Reserve válido elige un dueño; desde un estado reservado, la reserva de otro cliente se rechaza. Repetir el mismo \texttt{command\_id} devuelve el resultado previo sin duplicar la transición.

\par\smallskip\textbf{Criterio de autocorrección.} Debe declarar la transición inválida, la respuesta a un comando repetido y la identidad utilizada para reconocerlo.
\end{uaanswer}

\begin{uaexercise}
\textbf{6. Safety y liveness.} Formule dos propiedades de safety y dos condiciones de liveness para el cluster de cinco réplicas. Muestre una ejecución que conserva safety pero pierde liveness.
\end{uaexercise}
\begin{uaanswer}
Safety: no hay dos commits incompatibles por índice y VIP-01 no tiene dos dueños. Liveness: una solicitud válida progresa si hay mayoría comunicada y líder estable. Una partición permanente 3|2 preserva safety, mientras la minoría no progresa.

\par\smallskip\textbf{Criterio de autocorrección.} Debe distinguir una propiedad que nunca puede violarse de las condiciones bajo las cuales el sistema eventualmente progresa.
\end{uaanswer}

\begin{uaexercise}
\textbf{7. Raft y frontera de commit.} Complete los logs de N1--N5 cuando N1 propone la entrada 87 en term 12. Marque cuándo está local, replicada, committed, aplicada y confirmada.
\end{uaexercise}
\begin{uaanswer}
N1 agrega (87,12,Reserve Sofia); con N2 hay dos copias, aún no mayoría; con N3 hay 3/5 y N1 puede avanzar commitIndex. N1 aplica localmente y responde; followers aplican cuando reciben leaderCommit. Cada frontera debe marcarse por separado.

\par\smallskip\textbf{Criterio de autocorrección.} Debe marcar cinco fronteras distintas: append local, réplica, mayoría, aplicación y respuesta al cliente.
\end{uaanswer}

\begin{uaexercise}
\textbf{8. Crash móvil.} Analice dos crashes del líder: uno después del append local y otro después de la mayoría pero antes del ACK. Indique qué puede responder el nuevo líder y qué debería observar Sofía.
\end{uaexercise}
\begin{uaanswer}
Antes de mayoría, la entrada puede desaparecer y no debía existir ACK. Después de mayoría, la decisión committed sobrevive aunque el ACK se pierda; al reintentar, la identidad de operación permite devolver el mismo resultado.

\par\smallskip\textbf{Criterio de autocorrección.} Debe explicar la ambigüedad del ACK perdido y cómo la idempotencia permite consultar o repetir sin crear una segunda reserva.
\end{uaanswer}

\begin{uaexercise}
\textbf{9. Total-order broadcast.} Explique por qué una máquina de estados replicada necesita los mismos comandos en el mismo orden. Dé un contraejemplo donde timestamps Lamport no basten.
\end{uaexercise}
\begin{uaanswer}
Si réplicas aplican A,B y otras B,A, una máquina no conmutativa puede divergir. Lamport puede ordenar etiquetas, pero no garantiza que todas reciban el mismo conjunto ni evita huecos. Consenso repetido sí construye una secuencia comprometida común.

\par\smallskip\textbf{Criterio de autocorrección.} Debe mostrar una operación no conmutativa o un conjunto de mensajes distinto entre réplicas; ordenar etiquetas por sí solo no basta.
\end{uaanswer}

\begin{uaexercise}
\textbf{10. Pausa, GC y fencing.} N1 se pausa; N2 es elegido en un término mayor; N1 despierta y escribe sobre un recurso externo. Diseñe una protección mediante término y otra mediante fencing token.
\end{uaexercise}
\begin{uaanswer}
Raft rechaza mensajes de term 12 cuando el cluster ya está en 13. Para un recurso externo, el líder incluye token 13; el recurso recuerda el mayor token y rechaza 12. GC es una pausa del proceso, no una prueba de muerte.

\par\smallskip\textbf{Criterio de autocorrección.} Debe indicar quién emite el token, quién recuerda el máximo y qué operación concreta rechaza el recurso externo.
\end{uaanswer}

\begin{uaexercise}
\textbf{11. Raft y Multi-Paxos.} Mapee term, index, leader y followers a ballot, slot, proposer y acceptors para la decisión 87. Identifique la abstracción común y una diferencia operacional.
\end{uaexercise}
\begin{uaanswer}
Raft: term 12, index 87, líder y followers. Multi-Paxos: ballot 12, slot 87, proposer estable y acceptors. Ambos eligen un valor seguro por posición; difieren en organización, mensajes y reglas de log/liderazgo.

\par\smallskip\textbf{Criterio de autocorrección.} Debe mapear una misma posición y valor, sin afirmar que los protocolos tienen mensajes o estructura interna idénticos.
\end{uaanswer}

\begin{uaexercise}
\textbf{12. Tres logs.} Para log de consenso, WAL local y log Kafka, indique productor, unidad, garantía, recovery y una confusión frecuente.
\end{uaexercise}
\begin{uaanswer}
Consenso: comandos autoritativos y orden. WAL: bytes necesarios para recovery local. Kafka: eventos para consumo y replay. Confundirlos lleva a creer que publicar decide el estado o que commit distribuido elimina la persistencia local.

\par\smallskip\textbf{Criterio de autocorrección.} Debe mantener separadas autoridad global, recuperación local y distribución de eventos; cada log responde una pregunta diferente.
\end{uaanswer}

\begin{uaexercise}
\textbf{13. CAP sin slogan.} Durante una partición 3|2, Sofía ya fue committed en la mayoría y Tomás llega a la minoría. Compare responder éxito, rechazo y pendiente. Identifique qué propiedad se sacrifica.
\end{uaexercise}
\begin{uaanswer}
Éxito en la minoría conserva disponibilidad local pero puede violar el invariante. Rechazo/espera conserva una historia fuerte y sacrifica disponibilidad CAP. Pendiente cambia la semántica: acepta intención, pero no confirma la venta.

\par\smallskip\textbf{Criterio de autocorrección.} Debe indicar qué respuesta conserva el invariante y cuál mantiene disponibilidad, y aclarar que ``pendiente cambia la semántica.
\end{uaanswer}

\begin{uaexercise}
\textbf{14. PACELC y recencia.} Compare una lectura local de 15 ms posiblemente atrasada con una coordinada de 180 ms. Proponga dos operaciones que elijan políticas distintas y métricas para evaluarlas.
\end{uaexercise}
\begin{uaanswer}
Perfil o feed pueden usar lectura local con métrica de staleness; saldo o propiedad pueden coordinar. Medir p95/p99, antigüedad de dato, tasa de rechazo y disponibilidad semántica.

\par\smallskip\textbf{Criterio de autocorrección.} Debe incluir al menos una métrica de latencia y otra de frescura o staleness, no solo una preferencia cualitativa.
\end{uaanswer}

\begin{uaexercise}
\textbf{15. Consistencia.} Diferencie recencia, linearizabilidad, serializabilidad y consistencia externa mediante cuatro historias mínimas.
\end{uaexercise}
\begin{uaanswer}
Recencia mide frescura; linearizabilidad respeta tiempo real de operaciones no superpuestas; serializabilidad equivale a algún orden serial; consistencia externa combina orden serial con tiempo real observado por clientes.

\par\smallskip\textbf{Criterio de autocorrección.} Debe usar historias temporales concretas y no definir las cuatro garantías como sinónimos de ``dato actualizado.
\end{uaanswer}

\begin{uaexercise}
\textbf{16. Mayoría vs quorum.} Compare mayoría Raft con N=5,R=3,W=3 en un sistema Dynamo-like. Explique por qué tres respuestas no significan lo mismo.
\end{uaexercise}
\begin{uaanswer}
En Raft, la mayoría participa en la elección de una entrada única bajo términos y matching. En Dynamo, R/W completan operaciones sobre versiones de una clave. Puede haber hermanos concurrentes aun cuando R+W>N.

\par\smallskip\textbf{Criterio de autocorrección.} Debe explicar qué prueba aporta la intersección en cada modelo y qué reglas adicionales son necesarias para elegir una historia única.
\end{uaanswer}

\begin{uaexercise}
\textbf{17. Límites de R+W>N.} Construya tres contraejemplos: escrituras concurrentes, sloppy quorum y last-write-wins con reloj desalineado. En cada uno explique qué intersección existe y qué garantía falta.
\end{uaexercise}
\begin{uaanswer}
Concurrentes: el nodo común conserva dos ramas. Sloppy: el conjunto efectivo está fuera de propietarios. LWW: reloj adelantado descarta la rama válida. En los tres, la intersección no aporta el protocolo o semántica necesarios para linearizabilidad.

\par\smallskip\textbf{Criterio de autocorrección.} Debe construir cada contraejemplo con nodos y versiones concretos, señalando exactamente dónde deja de alcanzar la intersección.
\end{uaanswer}

\begin{uaexercise}
\textbf{18. Versiones y merge.} Modele un carrito con dos versiones concurrentes y diseñe un merge. Luego aplique la misma estrategia a VIP-01 y explique por qué falla.
\end{uaexercise}
\begin{uaanswer}
Carrito: unión de ítems con reglas para cantidades/remociones y nuevo vector que domina ambas ramas. VIP-01: unir propietarios rompe el invariante; requiere autoridad, consenso o rechazo, no merge.

\par\smallskip\textbf{Criterio de autocorrección.} Debe definir una política de merge determinista para el carrito y demostrar el invariante que esa política rompería en VIP-01.
\end{uaanswer}

\begin{uaexercise}
\textbf{19. Dynamo y reparación.} Diseñe una ejecución con nodo preferido caído, sloppy quorum, hint, hinted handoff, read repair y anti-entropy. Defina cuatro métricas de deuda de reparación.
\end{uaexercise}
\begin{uaanswer}
La escritura va a un nodo alternativo con hint para el propietario. Al regresar se hace handoff; una lectura puede reparar copias atrasadas; anti-entropy recorre rangos. Métricas: hints pendientes, edad máxima, bytes divergentes, repair/s, lag de reparación y conflictos por lectura.

\par\smallskip\textbf{Criterio de autocorrección.} Debe mostrar el ciclo completo desde la escritura desplazada hasta la reparación y proponer métricas de volumen, edad y capacidad.
\end{uaanswer}

\begin{uaexercise}
\textbf{20. Spanner y TrueTime.} Una transacción multi-grupo recibe $TT.now()=[4,16]$ ms y el coordinador elige $s=16$ ms. Muestre cuándo puede hacerse visible, explique external consistency y analice qué cambia si aumenta la incertidumbre.
\end{uaexercise}
\begin{uaanswer}
El coordinador elige $s=16$. No puede publicar mientras earliest sea 4, 8 o cualquier valor $\le16$; puede hacerlo cuando earliest>16, por ejemplo 17. Commit wait compra orden externo. Mayor incertidumbre ensancha el intervalo y normalmente aumenta la espera.

\par\smallskip\textbf{Criterio de autocorrección.} Debe identificar al coordinador, la elección de s, la desigualdad earliest>s y el efecto de ampliar la incertidumbre temporal.
\end{uaanswer}
\end{document}
